\documentclass[12pt]{article}
\usepackage[margin=1in]{geometry}
\usepackage{parskip}
\usepackage{titlesec}
\usepackage{enumitem}
\usepackage[most]{tcolorbox}
\tcbuselibrary{breakable}
\usepackage{hyperref}
\usepackage{bookmark}
\usepackage{setspace}
\usepackage{graphicx}
\usepackage{xcolor}

% Define colors for mechanics
\definecolor{mechblue}{RGB}{0,100,200}
\definecolor{mechgreen}{RGB}{0,150,0}
\definecolor{mechred}{RGB}{200,0,0}
\definecolor{mechgray}{RGB}{80,80,80}

% Mechanics box style
\newtcolorbox{mechanic}[1][]{
  colback=mechgray!5,
  colframe=mechgray!70!black,
  title=#1,
  fonttitle=\bfseries,
  left=2mm,
  right=2mm,
  top=2mm,
  bottom=2mm
}

% Quote style
\renewcommand{\quote}{
  \list{}{%
    \leftmargin=2em%
    \rightmargin=2em%
  }%
  \item[]%
}
\let\endquote=\endlist

% Title formatting
\titleformat{\section}
  {\normalfont\Large\bfseries}{\thesection}{1em}{}
\titleformat{\subsection}
  {\normalfont\large\bfseries}{\thesubsection}{1em}{}
\titleformat{\subsubsection}
  {\normalfont\normalsize\bfseries}{\thesubsubsection}{1em}{}

\begin{document}

\thispagestyle{empty}

\section*{The Ikari: Vilikari and Ostrikari}
\textit{``We are the first to plow this soil and the last to bury our dead here. The steppe tests us, the Ykrul hunger for us, and the Black Banners buy us---but the land remembers whose hands held the sickle.''}\\[1em]
\textit{---Old Vilikari proverb, carved into a boundary stone near the Amaranthine shores}

\subsection{Overview}
The Ikari---comprising the western Vilikari and eastern Ostrikari---are the indigenous peoples of the continent south of the Aelerian Mountains and north of the Amaranthine Sea. They inhabit the \textbf{Violet Steppe} (also called the Western Steppe or Prairieland), a vast expanse of tallgrass prairie and rolling hills that transitions from arid north to fertile river valleys near the sea. Unlike the nomadic Ykrul or the mercenary Black Banner companies, the Ikari are primarily agrarian and village-based, living in fortified \textit{burgs} surrounded by communal fields and pastures. Their society centers on the \textit{thing} (assembly), the \textit{hof} (extended family hall), and the sacred bond between farmer and field.

The Ikari exist in a constant three-way tension: 
\begin{itemize}[leftmargin=*]
  \item \textbf{With the Ykrul} to the north and east: Nomadic horse-lords who demand tribute, raid for slaves and horses, and view Ikari farms as rich pasture to be exploited.
  \item \textbf{With the Black Banner companies} along the Amaranthine coast and river valleys: Mercenary bands who sell their swords to the highest bidder---sometimes protecting Ikari settlements from Ykrul raids, sometimes conquering them for land and tax revenue.
  \item \textbf{With the land itself}: The Violet Steppe rewards patience and communal effort but punishes individualism; failed harvests or broken oaths invite famine, disease, or the creeping influence of forgotten spirits in the soil.
\end{itemize}

Despite pressures, the Ikari endure through deep-rooted traditions of oath-keeping, seasonal festivals, and a fierce belief that \textit{their} blood mingles with the steppe's soil---a belief that grants them uncanny resilience when defending their homesteads but carries a cost when they stray too far from home.

\subsection{Genre and Tone}
\textbf{Genre:} Low fantasy / migration period / frontier struggle  
\textbf{Key themes:} Communal obligation vs. individual glory, the weight of ancestral oaths, the fertility of blood-soil, the cost of settling a contested land  
\textbf{Dominant powers:} Ikari thing-assemblies (temporal), \textit{godi} (priest-chieftains, spiritual), wandering \textit{seidr-workers} (folk mystics, existential), the Violet Steppe itself (as a sentient, memory-laden entity)

\subsection{Attitude Toward Magic}
The Ikari practice a pragmatic, earth-bound magic woven into daily life---blessings for seeds, wards against livestock sickness, omens read in bird flights. They respect \textit{seidr} (folk mystics who commune with land-spirits) but distrust overt sorcery that severs the bond between people and soil. Visible spellcasting without agricultural or communal purpose is seen as \textit{unsteadying}---a sign the caster has forgotten their oath to the land. Licensed thaumaturgy (like the Aeler"s) is tolerated only when tied to specific, beneficial works (e.g., irrigation, pest control); otherwise, it invites suspicion. The Ikari say: \textit{"Magic that does not feed the field or heal the hof steals breath from the steppe."}

\subsection{Typical Dress and Customs}
\begin{itemize}[leftmargin=*]
  \item \textbf{Dress:} Undyed wool tunics and trousers, leather belts with iron buckles, cloaks fastened with thorn-shaped brooches. Warriors wear layered leather armor (\textit{lorica hamata} inspired) and conical helmets. Both genders braid hair; warriors often shave one side of the head. Footwear is soft leather turnshoes.
  \item \textbf{Customs:} 
  \begin{itemize}[leftmargin=*]
    \item Oaths sworn on soil or blade are unbreakable; breaking one invites \textit{nið} (spiritual shame) that draws misfortune.
    \item Harvest festivals (\textit{haustblót}) involve feasting, horse fights, and the ritual burial of the first sheaf to ensure next year's yield.
    \item Disputes are settled at the \textit{thing}; punishments range from fines (paid in livestock) to outlawry (\textit{skóggangr})---exile to the wild steppe where Ykrul hunt freely.
    \item Dead are buried in mound graves with personal goods; cremation is reserved for warriors fallen in defense of the burg.
  \end{itemize}
\end{itemize}

\subsection{Player Hook}
You were born in an Ikari burg where the soil remembers your ancestors" footsteps---but now the Ykrul press harder each spring, the Black Banners demand ever-larger tribute, and the old gods whisper through the wind that the thing-assembly has grown weak. Do you take up the spear to defend your hof, seek fortune as a mercenary with the Black Banners, or listen to the steppe's call to seek forgotten power in the burial mounds?

\subsection{Key Stakeholders}
\begin{itemize}[leftmargin=*]
  \item \textbf{The Godi:} Priest-chieftains who lead the thing, interpret omens, and maintain ties to land-spirits. Some walk the path of the \textit{seidr-worker}, communing directly with the Violet Steppe's memory.
  \item \textbf{The Warband Leaders:} Chieftains who command the \textit{leið} (local militia) and \textit{lið} (elite warrior retinue). Their loyalty is bought with weapons, horses, and promises of plunder.
  \item \textbf{The Free Farmers:} The backbone of Ikari society---those who own a \textit{boer} (farmstead) and owe service to the thing. They value peace but will fight to the death for their fields.
  \item \textbf{The Thralls:} Captives (often Ykrul or Black Banner deserters) bound to work the fields. Some earn freedom through valor; others remain bound until death.
\end{itemize}

\section{Lore Drops: The Ikari and the Violet Steppe}

\subsection{The Three Tensions}
The Ikari's existence is defined by pressures from all sides:
\begin{itemize}[leftmargin=*]
  \item \textbf{Ykrul Pressure:} The Ykrul see the Violet Steppe as their rightful hunting ground. Their raids peak at harvest time, targeting livestock and enslaving farmers. Some Ikari burgs pay \textit{skatt} (tribute in horses and grain) to buy peace---but the Ykrul always return for more.
  \item \textbf{Black Banner Influence:} Mercenary companies operate along the Amaranthine coast and major rivers. They offer protection contracts to Ikari burgs, but their loyalty shifts with coin. Worse, some companies have begun settling and claiming land as their own, displacing Ikari farmers.
  \item \textbf{Steppe's Memory:} The Violet Steppe is not inert soil---it remembers oaths sworn upon it and blood spilled upon it. Ikari who uphold their communal oaths find the earth yielding and their steps light; those who break troth or flee their duties suffer \textit{jordtrötthet} (soil-fatigue)---a creeping exhaustion that makes farming futile and draws whispers from the burial mounds.
\end{itemize}

\subsection{Life on the Steppe}
Ikari society balances settlement and vigilance:
\begin{itemize}[leftmargin=*]
  \item \textbf{Burg and Field:} Each burg centers on a fortified hall (\textit{hof}) surrounded by wooden palisades. Fields are divided into \textit{teilar} (strips) reallocated yearly at the thing to ensure fairness. Pastures are held communally, with herds moved seasonally.
  \item \textbf{The Thing-Assembly:} Free farmers gather at the sacred \textit{thingstead} to settle disputes, allocate land, and choose warleaders. Decisions require consensus; dissenters may be challenged to \textit{hólmganga} (duel) to prove their stance.
  \item \textbf{Seasonal Rhythms:} 
  \begin{itemize}[leftmargin=*]
    \item \textit{Spring:} Ploughing, oath-renewal at the thing, vigilance against early Ykrul probes.
    \item \textit{Summer:} Tending fields, horse breeding, preparation for harvest.
    \item \textit{Autumn:} Harvest, \textit{haustblót} festival, tribute payment to Ykrul or Black Banners (if forced).
    \item \textit{Winter:} Indoor crafts, storytelling, planning for next year; slowest raid season but highest risk of famine.
  \end{itemize}
  \item \textbf{The Wandering Seidr:} Mystics who leave the burg to commune with burial mounds, springs, and standing stones. They bring back omens, healing knowledge, and sometimes dangerous bargains with the steppe's deeper memories.
\end{itemize}

\subsection{Relations with Neighbors}
\begin{itemize}[leftmargin=*]
  \item \textbf{With the Aeler:} Limited but tense. Aeler vent-keeps dot the eastern foothills, trading salt and iron for Ikari grain and wool. The Aeler view the Ikari as useful but undisciplined; the Ikari see the Aeler as breath-counters who miss the steppe's soul.
  \item \textbf{With the Linn:} Minimal direct contact (separated by the Direwood and Mistlands), but Ikari traders sometimes acquire Linn amber or ironwork via intermediaries. Shared reverence for ancestral oaths creates subtle kinship.
  \item \textbf{With the Ykrul:} Beyond raids, there is periodic trade---Ikari grain for Ykrul furs and horses---and occasional alliances against common threats (e.g., Black Banner overreach).
  \item \textbf{With the Black Banner Companies:} A relationship of necessity. Ikari hire them as \textit{landsvern} (land-defenders) against Ykrul, but remain wary of their long-term goals. Some Ikari chieftains seek to become Black Banner captains themselves.
\end{itemize}

\section{Mechanical Breakouts}

\subsection{Cultural Talents}
The Ikari possess unique abilities rooted in their bond to the Violet Steppe and communal oath-keeping.

\begin{mechanic}{Minor Talents (2--3 XP)}
  \begin{itemize}[leftmargin=*]
    \item \textbf{Steppe-Watcher [PERCEPTION]:} Passive. +1 die to notice changes in terrain, weather, or livestock behavior on the Violet Steppe. Grants advantage when predicting Ykrul raid patterns or Black Banner movements.
    \item \textbf{Oath-Keeper's Resolve [RESOLVE]:} Passive. When you swear an oath on soil or blade (requiring 1 minute of roleplay), gain +1 die to resist fear, charm, or coercion related to breaking that oath. If you break the oath, suffer --2 dice to all rolls until you make amends (as judged by the GM).
    \item \textbf{Burg-Sense [SURVIVAL]:} Passive. While within 1 mile of an Ikari burg (or other settled community you consider home), reduce Fatigue from lack of food or water by 1 level (minimum 0). Outside this range, suffer +1 Fatigue from exposure to the open steppe.
  \end{itemize}
\end{mechanic}

\begin{mechanic}{Major Talents (4--6 XP)}
  \begin{itemize}[leftmargin=*]
    \item \textbf{Thing-Speaker [Sway]:} Active, once per scene. When addressing a group involved in a dispute or decision (e.g., at a thing-assembly, in a warband, during negotiations), you may declare your intent to speak for the common good. If successful, gain +2 dice to Persuade or Command rolls, and listeners suffer --1 die to resist your argument unless they provide a compelling counter-oath.
    \item \textbf{Warband-Leader [Command]:} Passive. When leading a group of at least three allies in combat or a hazardous task, your Position improves by 1 step (Desperate → Controlled → Dominant). Requires: Command 2+.
    \item \textbf{Soil-Bound [Endurance]:} Passive. You suffer no Fatigue from marching or labor on the Violet Steppe. However, if you spend more than 3 consecutive days north of the Aelerian Mountains (in the lands of the Aeler or Ykrul deep steppe) or south of the Amaranthine Sea (in coastal or jungle lands), you begin to accumulate \textit{jordtrötthet}: each day beyond the third adds 1 Fatigue that cannot be removed by rest until you return to the steppe.
  \end{itemize}
\end{mechanic}

\begin{mechanic}{Prestige Talents (7--10 XP)}
  \begin{itemize}[leftmargin=*]
    \item \textbf{Gods-Marked [Spirit]:} Active, once per session. After spending 10 minutes in quiet communion at a burial mound, standing stone, or sacred spring (requiring a suitable location), you may invoke the steppe's memory. Gain +2 dice to one roll related to healing, divination, or protecting your burg. If used to harm the steppe or break a major oath, the GM gains 2 Story Beats.
    \item \textbf{Leið Commander [Tactics]:} Active, once per scene. When organizing a defensive action (e.g., holding a burg against siege, setting an ambush), you may declare your intent to defend the communal land. If successful, all allies gain +1 die to defend rolls, and enemies suffer --1 die to approach rolls. Requires: Tactics 3+, Command 2+.
  \end{itemize}
\end{mechanic}

\subsection{Signature Background: Ikari Burg-Born}
\textit{Access Tags:} Burg-Law, Thing-Stead Right  
\textit{Skills:} Survival 2, Athletics 1  
\textit{Signature Contact:} Godi of Your Burg (Cap 1 follower; will not take independent actions but grants +1 die once per scene when your actions align with burg welfare)  
\textit{Background Boon:} Once per session when acting to defend your burg, its fields, or its thing-assembly, treat one Position step as better (e.g., Desperate → Controlled) or reduce DV by 1 for a task tied to communal welfare.  
\textit{Obligation Seed:} The Burg's Oath (4 segments). When filled, the thing-assembly demands you answer for neglecting your duties---through public penance, a dangerous task, or payment of fines in livestock. Cleared by fulfilling the burg's needs (e.g., repairing palisades, harvesting communal fields, standing watch).

\subsection{Terrestrial Patron Option: The Thing-Stead}
Some Ikari swear oaths not to a burg, but to the \textit{thingstead} itself---the sacred assembly place where oaths are made and laws spoken. This functions as a terrestrial patron with:
\begin{itemize}[leftmargin=*]
  \item \textbf{Perks:} Sanctuary at any thing‑stead (safe haven, counsel from elders), advantage on rolls to resolve disputes peacefully, access to thing‑stead‑specific omens (via Lore + Spirit rolls).
  \item \textbf{Demands:} Attend the thing‑stead at least once per season; uphold all thing‑decisions even if personally disagreeing; never draw steel within the thing‑stead's bounds except in sanctioned \textit{hólmganga}.
  \item \textbf{Obligation:} Tracks as standard (Spirit + Presence). When full, the thing‑stead's spirit withdraws its favor---you suffer --2 dice to all rolls involving leadership or oath‑keeping until you make a public act of atonement (e.g., confessing fault at the thing‑stead, rebuilding a damaged marker stone).
\end{itemize}
\noindent\textit{Example:} A character with the Thing‑Stead patron might gain the boon \textit{Thing‑Speaker's Voice} (once per session, +2 dice to Persuade when arguing for the thing‑stead's decision) but risks obligation if they skip a thing‑meeting to pursue personal glory.

\subsection{Tags and Mechanics}
\begin{itemize}[leftmargin=*]
  \item \textbf{[STEPPE-WARDEN]:} Granted by talents like Soil-Bound or Burg-Sense. When active, treats the Violet Steppe as favorable terrain (no movement penalties, advantage on Stealth/Perception in tallgrass).
  \item \textbf{[HERD-BOUND]:} Granted by Burg-Sense when near livestock. Allows spending 1 Boon to ignore Fatigue from herding or defending animals against predators.
  \item \textbf{[OATH-BOUND]:} Granted by Oath‑Keeper's Resolve. When active, breaking a sworn oath triggers immediate GM intervention (e.g., a Story Beat for misfortune, a Complication for spiritual backlash).
  \item \textbf{[THING-SPEAKER]:} Granted by Thing‑Speaker talent. When active, allows spending 1 Boon to force a group to hear you out in a dispute (no roll required, but listeners may still reject your argument after hearing it).
\end{itemize}

\subsection{Integration with Core Mechanics}
\begin{itemize}[leftmargin=*]
  \item \textbf{Position:} Ikari characters gain +1 die to Position-related rolls (e.g., improving Position, resisting Position worseners) when acting to defend their burg, fields, or thing‑stead---reflecting the steppe's favor for those who uphold communal oaths. Conversely, they suffer --1 die when acting purely for personal gain against the burg's interests (e.g., stealing from the communal store, abandoning a watch post).
  \item \textbf{Boons:} Ikari characters often earn Boons through acts of communal defense (e.g., repelling a Ykrul raid, resolving a thing‑stead dispute) or upholding difficult oaths. Spending Boons to improve Position or activate assets (like a burg's defenses) is especially narratively fitting.
  \item \textbf{Obligation:} Ikari characters frequently accrue obligation not just from patrons, but from the burg or thing‑stead itself---representing the weight of communal duty. Clearing this obligation through service (e.g., rebuilding after a raid, fasting for the burg's welfare) is a core narrative loop.
\end{itemize}

\section{Adventure Seeds}

\subsection{1. The Burg's Oath}
A neighboring Ikari burg has failed to pay its tribute to the Ykrul khan---and now the khan demands not just tribute, but the burg's godi as hostage. The thing‑stead has decreed that the burg must be abandoned to save the region, but you believe the godi can be rescued. Do you defy the thing‑stead to save your kin, or uphold the oath to the collective and watch your burg fall?

\subsection{2. The Black Banner's Bargain**
A Black Banner captain offers your burg a deal: they will build a stone palisade and train your levy in exchange for half your harvest for three years. The thing‑stead is split---some see salvation, others see slow conquest. Investigating further, you discover the captain plans to claim the burg's land as his own after the term ends. Do you warn the thing‑stead (risking being called a traitor for meddling), take the deal and plan a counter‑move, or refuse and prepare for immediate Ykrul pressure?

\subsection{3. The Seidr's Whisper
You"ve been having dreams of a burial mound crying out in the night---dreams that leave you with dirt under your nails and a taste of iron. A local seidr-worker interprets this as the steppe itself calling for a forgotten oath to be renewed, possibly tied to a ancient burg swallowed by the steppe centuries ago. Following the dream leads you deep into Ykrul‑claimed territory, where the mound is guarded not by men, but by something that walks on too‑many legs and speaks in the voice of your ancestors. Do you renew the oath (risking what the steppe demands in return), or seal the mound forever to silence its call?

\subsection{4. The Three‑Way Truce (Campaign Arc)
A rare moment has arrived: the Ykrul khan, the Black Banner marshal, and the Ikari thing‑stead have agreed to a summit at a neutral river ford to discuss boundaries and tribute. You are chosen as part of the Ikari delegation---but old grudges run deep, and all three sides bring hidden agendas. Success means a generation of peace; failure means renewed raids, mercenary overreach, or worse---something older than all three stirring in the steppe's roots as the oaths of men are broken.

\end{document}
